\begin{block}{Why Kugire?}
\begin{itemize}
\item We perform analysis on interpretive diversity for visible lexical elements in waka
\item Analysis on invisible discourse-level elements was not performed yet
\item Kugire shows discourse-level interpretability
\end{itemize}

\vspace{-\dimexpr\parskip+\topsep+\partopsep+\parsep+.65\baselineskip\relax}
{\setlength{\intextsep}{0pt}%
\setlength{\textfloatsep}{0pt}%
\setlength{\floatsep}{0pt}%
\begin{figure}[H]
\begin{minipage}[t]{0.55\linewidth}
  \centering
  \includesvg[inkscapelatex=false, width=0.9\textwidth]{../images/fig-miru}
  \setcounter{figure}{1}
  \caption{Interpretive diversity of the predicate of \textit{miru} (see)}
\end{minipage}%
\hfill
\begin{minipage}[t]{0.42\linewidth}
  \centering
  \includesvg[inkscapelatex=false, width=0.9\textwidth]{../images/fig-ominaeshi}
  \setcounter{figure}{2}
  \caption{Interpretive diversity of lexical space \textit{ominaeshi} (``maidenflower'')}
\end{minipage}
\end{figure}
}
\end{block}
